% --- Introduzione ---
\chapter{Introduzione}

Negli ultimi anni, la crescente interconnessione dei sistemi industriali e dei Supervisory Control and Data Acquisition (SCADA) ha esposto le infrastrutture critiche a minacce informatiche sempre più sofisticate.  
La ricerca ha dimostrato l’efficacia di approcci multilivello e semantici per la protezione dei sistemi ICS, l’utilizzo di tecniche di Machine Learning per l’identificazione delle anomalie e la simulazione di scenari realistici per la validazione delle difese esistenti \cite{simulationImplementationSecurityServices, diderotIds, semanticSecurityAnalysisScadaNetwork, graphBasedActiveCyberDefense}.

Il progetto S.I.D.E. nasce dall’esigenza di disporre di un ambiente controllato che permetta di testare strumenti di difesa, metodologie di rilevazione delle intrusioni e strategie di monitoraggio avanzate.  
A tal fine, S.I.D.E. integra simulatori di protocolli industriali (Modbus TCP, Siemens S7 e OPC-UA), consentendo la creazione di traffico client-server realistico, l’inserimento di anomalie programmate e la raccolta di dati dettagliati per analisi successive.

L’obiettivo principale è fornire un framework modulare e scalabile, utile per l'ambito industriale, che permetta di:
\begin{itemize}
    \item simulare scenari ICS/SCADA realistici e variabili;
    \item valutare l’efficacia di sistemi di rilevazione delle intrusioni e meccanismi di difesa attiva;
    \item analizzare il comportamento dei protocolli industriali e delle architetture SCADA in condizioni normali e anomale;
\end{itemize}

\section{Struttura del documento}
Il presente documento è articolato in diversi capitoli, ognuno dei quali approfondisce aspetti specifici del lavoro svolto.

\begin{itemize}
    \item \textbf{Capitolo 2} introduce e analizza i principali lavori correlati riguardanti la sicurezza dei sistemi \textit{ICS/SCADA}, con particolare attenzione ai modelli di difesa basati su \textit{Intrusion Detection Systems (IDS)}, approcci di \textit{machine learning}, analisi semantica e metodologie di difesa attiva che rappresentano l’evoluzione delle strategie di protezione delle infrastrutture critiche.    
    \item \textbf{Capitolo 3} descrive in dettaglio l’organizzazione del progetto, la struttura architetturale, i vari \textit{sub-moduli} sviluppati e le tecnologie utilizzate, con particolare attenzione ai protocolli SCADA oggetto di studio.  
    \item \textbf{Capitolo 4} riassume le conclusioni emerse dal lavoro, proponendo possibili miglioramenti e sviluppi futuri del sistema.  
\end{itemize}
